\section{长上下文、记忆与持续学习：能力边界正在重新划线}

\subsection{Context Length 进步与“压缩焦虑”}

上下文窗口在持续增长，但访谈强调：长上下文不是免费午餐。窗口变大后，如何检索关键信息、如何压缩不丢语义、如何在长链路里保持一致性，才是新难点。很多工具中的“自动压缩”会把高价值细节折叠成粗糙摘要，导致后续推理失真。

\begin{figure}[H]
\centering
\includegraphics[width=0.88\textwidth]{frame-03.jpg}
\caption{访谈中关于长上下文与压缩策略的讨论\protect\footnotemark}
\end{figure}
\footnotetext{视频画面时间区间：02:44:10--02:44:20。}

\begin{importantbox}{长上下文系统的三个工程关键}
\begin{itemize}[leftmargin=1.5em]
    \item \textbf{检索策略}：不是“全喂给模型”，而是先筛后读；
    \item \textbf{压缩策略}：保存可执行事实，丢弃低价值修辞；
    \item \textbf{恢复策略}：摘要失真时能回溯原文片段并纠偏。
\end{itemize}
\end{importantbox}

\subsection{Continual Learning：权重更新还是上下文注入}

访谈把持续学习拆成两条路：更新权重（真正学习）和注入记忆（运行时补充）。前者成本高、风险高，但长期更稳；后者实现快、成本低，但易受上下文窗口与检索质量限制。现实系统常采用混合路线。

\begin{knowledgebox}{LoRA 与在线更新的现实权衡}
LoRA 这类低秩适配方法在企业定制中常见，因为它在“学习速度”和“遗忘风险”之间给出可控折中。若追求分钟级在线更新，则需强监控和回滚能力，否则很容易把噪声反馈写进模型行为。
\end{knowledgebox}

\begin{warningbox}{“记住用户”并不等于“理解用户”}
把用户历史堆进上下文可以改善个性化，但也可能放大旧偏差、引入隐私泄露风险，并让模型在新任务上过度依赖旧模板。个性化系统必须有明确的记忆边界与删除机制。
\end{warningbox}

\subsection{本章小结}

长上下文和持续学习是能力上限的关键变量，但真正难题在工程控制：检索、压缩、回滚、隐私与评价必须同时设计。

